\documentclass[french, 11pt]{article}

\input{setup.tex}

\def\chapitre{5}
\def\pagetitle{Stratégies Algorithmiques}

\begin{document}

\input{title.tex}

\section{Stratégies.}

\begin{defi}{Force brute.}{}
    \bf{Idée.} Explorer tous les candidats possibles.\\
    \bf{Contraintes.} Univers dénombrable, on sait tester si une entrée est solution.
\end{defi}

\vspace*{-0.7cm}

\begin{defi}{Backtracking.}{}
    \bf{Idée.} Construire une solution pas à pas et revenir sur le dernier choix en cas d'impasse.
\end{defi}

\vspace*{-0.7cm}

\begin{defi}{Algorithmes gloutons.}{}
    \bf{Idée.} Parier sur les maxima locaux.\\
    \bf{Iontraintes.} La solution n'est pas toujours optimale.
\end{defi}

\vspace*{-0.7cm}

\begin{defi}{Programmation dynamique.}{}
    Cette méthode est envisageable si :\\
    \bf{1. Sous-problèmes optimaux:} La solution pour une entrée donnée s'exprime en fonction des solutions pour des entrées strictement plus petites.\\
    \bf{2. Chevauchement de sous-problèmes:} La solution naïve mène à calculer plusieurs fois les mêmes solutions.\\
    On peut l'implémenter de deux manières :\\
    \bf{Approche de bas en haut:} Calculer les résultats dans l'ordre pour avoir les solutions quand on en a besoin. On peut les garder en mémoire (tableau...)\\
    \bf{Mémoïsation:} On crée une structure stockant les solutions. Au moment du calcul, on vérifie si la solution a déjà été calculée, sinon on la rajoute.
\end{defi}

\section{Exemples.}

\begin{ex}{Rendu de monnaie.}{}
    \bf{Problème.} On dispose d'un nombre illimité de pièces de valeurs $n_1>n_2>...>n_k$. Comment arriver à une certaine somme $S$ avec le moins de pièces ?\n
    L'algorithme glouton naturel consiste à puiser dans les pièces par ordre décroissant de valeurs sans dépasser $S$.
\end{ex}

\begin{ex}{Placement d'activités.}{}
    \bf{Problème.} attribuer des salles pour le plus de cours possibles.\\
    \bf{Théorème.} Le choix du cours terminant le plus tôt est optimal.
    \tcblower
    Supposons que ce choix mène à la liste de cours $(c_1,...c_n)$. Montrons par l'absurde qu'on ne peut pas avoir de suite plus longue par récurrence descendante sur la longueur du plus long préfixe commun.\\
    \bf{Initialisation.} $(d_1,...,d_n,...)$ avec $d_i=c_i$ pour $i\in\lb1,n\rb$.\\
    On peut alors construire $(c_1,...,c_n,c_{n+1})$, car $d_{n+1}$ est plaçable : contradiction.\\
    \bf{Hérédité.} Pour $k<n$, il ne peut pas exister une suite de cours plaçable de longueur $n+1$ ayant un préfixe commun de longueur au moins $k+1$ avec $(c_1,...,c_n)$.\\
    Supposons qu'il existe une suite de cours plaçables $(c_1,...,c_k,d_{k+1},...,d_{n+1})$, $c_{k+1}$ est plaçable entre $c_k$ et $d_{k+1}$, ce qui est absurde.
\end{ex}

\vspace*{-0.7cm}

\begin{ex}{Distance de Levenshtein}{}
    C'est une distance sur l'ensemble des mots d'un dictionnaire. Elle représente le nombre minimal d'opérations élémentaires pour passer d'un mot à un autre : substitution, délétion et insertion.
    Soient deux mots $u$ et $v$. On note $d[i,j]$ la distance de Levenshtein entre le préfixe de longueur $i$ de $u$ et le préfixe de longueur $j$ de $v$.
    \begin{equation*}
        d[i,j] = \min\begin{cases}
            d[i,j-1] + 1\\
            d[i-1,j] + 1\\
            d[i-1,j-1] + \delta_{u[i] \neq v[j]}
        \end{cases}
    \end{equation*}
    De plus, $d[0,j]=j$ et $d[i,0]=i$.\\
    \textbf{Complexité de l'algorithme naïf:}\\
    $T(m,n)=\alpha + T(m-1,n) + T(m,n-1) + T(m-1,n-1)$. En simplifiant: $T(m,n) \geq \alpha + 3T(m-1, n-1)$.\\
    On pose $U(n)=T(n,n)\geq\alpha+3U(n-1)$ donc $U(n)\geq\alpha n + 3^nU(0)=\Omega(3^n)$.\\
    \textbf{Complexité de l'algorithme dynamique:}\\
    On fait un nombre constant d'opérations pour chaque axe du tableau donc une complexité en $O(mn)$.
\end{ex}

\vspace*{-0.7cm}

\begin{ex}{Typographie}{}
    \textbf{Hypothèses:} Police de longueur fixe, un seul espace entre deux mots.\\
    \textbf{Entrée:} Une suite de $n$ mots de longueurs $l_1,...,l_n$ et la largeur $M$ de la ligne tels que $\forall i, ~ l_i \leq M$.\\
    \textbf{Contraintes:}\\
    $\hra$ Aucune ligne ne dépasse : si une ligne contient les mots $i$ à $j$, on a : $\sum_{k=i}^j\left(l_k - i + j\right) \leq M$.\\
    $\hra$ On minimise la somme cubes des espaces finaux sur chaque ligne, sauf la dernière.\\
    Montrons que la propriété de sous-problèmes optimaux est vérifiée.\\
    Soit une solution optimale sur $h$ lignes, alors la composition est optimale sur les $h-1$ premières lignes.\\
    Notons $\m{E}(i)$ la somme des cubes des espaces finaux de lignes, sauf la dernière à partir du mot $i$ jusqu'au dernier, en supposant que le mot $i$ est en début de ligne.\\
    $\hra$ Si on peut écrire tous les mots restants sur une seule ligne, il suffit de le faire $\E(i) = 0$.\\
    $\hra$ Sinon, il faut trouver $j$ tel que :
    \begin{equation*}
        \E(i)=\min\left\{\left(M-\left(\sum_{k=i}^jl_k+j-i\right)\right)^3+E(j+1) ~ | ~ j \in \llbracket1,n\rrbracket ~ \land ~ \sum_{k=i}^jl_k+j-i\leq M\right\}.
    \end{equation*}
\end{ex}

\fin

\end{document}